\documentclass{beamer}
\usetheme{Madrid}
\usecolortheme{default}

% Define custom colors for AI slides
\definecolor{AIBlue}{RGB}{30, 144, 255}
\definecolor{AILightBlue}{RGB}{173, 216, 230}
\definecolor{AINavy}{RGB}{25, 25, 112}

% Commands to switch to AI color scheme
\newcommand{\AIcolors}{%
  \setbeamercolor{frametitle}{bg=AIBlue,fg=white}%
  \setbeamercolor{title}{bg=AINavy,fg=white}%
  \setbeamercolor{structure}{fg=AIBlue}%
  \setbeamercolor{item}{fg=AINavy}%
  \setbeamercolor{normal text}{fg=black}%
}

% Command to restore default colors
\newcommand{\defaultcolors}{%
  \setbeamercolor{frametitle}{bg=osaorange,fg=white}%
  \setbeamercolor{title}{bg=osaorange,fg=white}%
  \setbeamercolor{structure}{fg=osaorange}%
  \setbeamercolor{item}{fg=black}%
  \setbeamercolor{normal text}{fg=black}%
}

% Define OSU orange color
\definecolor{osaorange}{RGB}{220,115,38}

% Set colors
\setbeamercolor{structure}{fg=osaorange}
\setbeamercolor{title}{fg=white,bg=osaorange}
\setbeamercolor{frametitle}{fg=white,bg=osaorange}

\usepackage{graphicx}
\usepackage{amsmath}
\usepackage{booktabs}
\usepackage{hyperref}
\usepackage{tikz}
\usetikzlibrary{shapes,arrows}

\title{H524 Introduction to Biostatistics}
\subtitle{Week 4: Confidence Intervals}
\author{Dr. John Molitor}
\institute{}
\date{Fall 2025}

\begin{document}

\frame{\titlepage}

\begin{frame}
\frametitle{Outline}
\tableofcontents
\end{frame}

\section{Review from Week 3}

\begin{frame}
\frametitle{Quick Review: Sampling Distributions}
\textbf{Last week we covered:}
\begin{itemize}
\item Binomial and normal distributions
\item Standardization and Z-scores
\item Central Limit Theorem: $\bar{X} \sim N\left(\mu, \frac{\sigma^2}{n}\right)$ for large $n$
\item Standard error: $\text{SE} = \frac{\sigma}{\sqrt{n}}$
\item Introduction to confidence intervals
\end{itemize}

\vspace{0.5cm}
\textbf{This Week:} Deep dive into confidence intervals (Chapter 9)

\vspace{0.3cm}
\textbf{Reading:} Pagano \& Gauvreau Chapter 9
\end{frame}

\begin{frame}
\frametitle{The Problem of Point Estimation}
\textbf{Scenario:} We measure cholesterol in 25 randomly selected adults
\begin{itemize}
\item Sample mean: $\bar{x} = 205$ mg/dL
\item Sample SD: $s = 42$ mg/dL
\end{itemize}

\vspace{0.5cm}
\textbf{Question:} What is the true population mean $\mu$?

\vspace{0.3cm}
\textbf{Point estimate:} Our best guess is $\bar{x} = 205$ mg/dL

\vspace{0.5cm}
\textbf{But...}
\begin{itemize}
\item Different samples would give different $\bar{x}$ values
\item We know there's sampling variability
\item How confident are we in this estimate?
\item What's the range of plausible values for $\mu$?
\end{itemize}

\vspace{0.3cm}
\textbf{Solution:} Use a \textit{confidence interval} instead of just a point estimate
\end{frame}

\section{Confidence Intervals: Core Concepts}

\begin{frame}
\frametitle{What is a Confidence Interval?}
\textbf{Definition:} A range of values that likely contains the unknown population parameter

\vspace{0.5cm}
\textbf{General Form:}
$$\text{Estimate} \pm \text{Margin of Error}$$
$$\text{Estimate} \pm (\text{Critical Value}) \times (\text{Standard Error})$$

\vspace{0.5cm}
\textbf{For a population mean:}
$$\bar{x} \pm t_{\alpha/2, n-1} \times \frac{s}{\sqrt{n}}$$

\vspace{0.5cm}
\textbf{Components:}
\begin{itemize}
\item $\bar{x}$ = sample mean (point estimate)
\item $s$ = sample standard deviation
\item $n$ = sample size
\item $t_{\alpha/2, n-1}$ = critical value from t-distribution
\end{itemize}
\end{frame}

\begin{frame}
\frametitle{Confidence Level and Interpretation}
\textbf{Confidence level:} Usually 95\%, but can be 90\%, 99\%, etc.

\vspace{0.5cm}
\textbf{Correct interpretation of a 95\% CI:}
\begin{itemize}
\item ``We are 95\% confident that the true population mean is between [lower, upper]''
\item If we repeated this procedure many times, 95\% of the intervals would contain $\mu$
\item The method is reliable 95\% of the time
\end{itemize}

\vspace{0.5cm}
\textbf{What 95\% confidence does NOT mean:}
\begin{itemize}
\item $\times$ ``There's a 95\% probability $\mu$ is in this interval''
  \begin{itemize}
  \item $\mu$ is fixed (not random); it's either in or out
  \end{itemize}
\item $\times$ ``95\% of the data falls in this interval''
  \begin{itemize}
  \item CI is about the mean, not individual values
  \end{itemize}
\item $\times$ ``There's a 95\% chance the next sample mean will be in this interval''
  \begin{itemize}
  \item Each new sample gets its own CI
  \end{itemize}
\end{itemize}
\end{frame}

\begin{frame}
\frametitle{Visualizing Confidence Intervals}
\begin{center}
\begin{tikzpicture}[scale=0.8]
% True mean line
\draw[thick] (0,6) -- (10,6);
\node at (10.5,6) {$\mu$ = 200};

% Multiple CIs
\foreach \y/\lower/\upper/\color in {
  5/3.5/6.8/green!60!black,
  4.5/3.2/6.5/green!60!black,
  4/3.8/6.2/green!60!black,
  3.5/4.2/6.9/green!60!black,
  3/2.5/5.8/red,
  2.5/3.9/6.4/green!60!black,
  2/3.6/6.7/green!60!black,
  1.5/4.1/6.3/green!60!black,
  1/3.4/6.1/green!60!black,
  0.5/3.7/6.6/green!60!black
} {
  \draw[\color, very thick] (\lower,\y) -- (\upper,\y);
  \filldraw[\color] ({(\lower+\upper)/2},\y) circle (2pt);
}

% Labels
\node[anchor=east] at (0,6) {True mean:};
\node[anchor=west, green!60!black] at (10.5,1) {Contains $\mu$};
\node[anchor=west, red] at (10.5,3) {Misses $\mu$};
\end{tikzpicture}
\end{center}

\vspace{0.3cm}
\textbf{If we took 20 samples:}
\begin{itemize}
\item Each sample gives a different CI
\item About 95\% of CIs (19 out of 20) contain $\mu$
\item About 5\% (1 out of 20) miss $\mu$
\end{itemize}
\end{frame}

\section{One-Sample t-Interval}

\begin{frame}
\frametitle{When to Use the t-Distribution}
\textbf{Recall from Week 3:}
\begin{itemize}
\item If $\sigma$ is known: use Z-distribution (rare in practice)
\item If $\sigma$ is unknown: use t-distribution (almost always)
\end{itemize}

\vspace{0.5cm}
\textbf{The t-distribution:}
\begin{itemize}
\item Bell-shaped like normal, but heavier tails
\item Accounts for uncertainty in estimating $\sigma$ with $s$
\item Determined by degrees of freedom: $\text{df} = n - 1$
\item As $n$ increases, t approaches normal
\end{itemize}

\vspace{0.5cm}
\textbf{Assumptions for t-interval:}
\begin{enumerate}
\item Random sample from the population
\item Population is approximately normal, OR $n$ is large ($n \geq 30$)
\item Observations are independent
\end{enumerate}

\vspace{0.3cm}
\textbf{Robust to mild departures from normality, especially for larger $n$}
\end{frame}

\begin{frame}
\frametitle{Computing a t-Based Confidence Interval}
\textbf{Step 1: Calculate sample statistics}
\begin{itemize}
\item Sample mean: $\bar{x}$
\item Sample SD: $s$
\item Sample size: $n$
\end{itemize}

\vspace{0.3cm}
\textbf{Step 2: Compute standard error}
$$\text{SE} = \frac{s}{\sqrt{n}}$$

\vspace{0.3cm}
\textbf{Step 3: Find critical value}
\begin{itemize}
\item Degrees of freedom: $\text{df} = n - 1$
\item For 95\% CI: $t_{0.975, \text{df}}$ (leaves 2.5\% in each tail)
\item Use t-table or R: \texttt{qt(0.975, df)}
\end{itemize}

\vspace{0.3cm}
\textbf{Step 4: Calculate margin of error}
$$\text{ME} = t_{\alpha/2, \text{df}} \times \text{SE}$$

\vspace{0.3cm}
\textbf{Step 5: Construct interval}
$$\text{CI} = \bar{x} \pm \text{ME} = (\bar{x} - \text{ME}, \bar{x} + \text{ME})$$
\end{frame}

\begin{frame}
\frametitle{Example: Blood Pressure Study}
\textbf{Research question:} What is the mean systolic blood pressure for adults with hypertension?

\vspace{0.3cm}
\textbf{Data:} Random sample of $n = 25$ hypertensive adults
\begin{itemize}
\item Sample mean: $\bar{x} = 128.4$ mmHg
\item Sample SD: $s = 12.6$ mmHg
\end{itemize}

\vspace{0.3cm}
\textbf{Calculate 95\% CI:}

\textbf{Step 1:} Standard error
$$\text{SE} = \frac{12.6}{\sqrt{25}} = \frac{12.6}{5} = 2.52 \text{ mmHg}$$

\textbf{Step 2:} Critical value ($\text{df} = 24$)
$$t_{0.975, 24} = 2.064$$

\textbf{Step 3:} Margin of error
$$\text{ME} = 2.064 \times 2.52 = 5.20 \text{ mmHg}$$

\textbf{Step 4:} Confidence interval
$$\text{95\% CI} = 128.4 \pm 5.20 = (123.2, 133.6) \text{ mmHg}$$
\end{frame}

\begin{frame}
\frametitle{Interpreting the Blood Pressure CI}
\textbf{Result:} 95\% CI = (123.2, 133.6) mmHg

\vspace{0.5cm}
\textbf{Correct interpretation:}
\begin{itemize}
\item ``We are 95\% confident that the true mean systolic blood pressure for adults with hypertension is between 123.2 and 133.6 mmHg''
\item Based on this sample, plausible values for $\mu$ range from 123.2 to 133.6
\item If we repeated this study many times, 95\% of our CIs would contain $\mu$
\end{itemize}

\vspace{0.5cm}
\textbf{Clinical context:}
\begin{itemize}
\item Normal BP: $< 120$ mmHg systolic
\item Elevated: 120-129 mmHg
\item Hypertension Stage 1: 130-139 mmHg
\item The entire CI is above normal range
\item This provides evidence that the population mean is elevated
\end{itemize}
\end{frame}

\begin{frame}
\frametitle{Example: Drug Efficacy Trial}
\textbf{Scenario:} A new pain medication is tested on 16 patients

\vspace{0.3cm}
\textbf{Outcome:} Pain reduction score (0-100 scale)
\begin{itemize}
\item Sample mean: $\bar{x} = 45.2$ points
\item Sample SD: $s = 8.4$ points
\item Sample size: $n = 16$
\end{itemize}

\vspace{0.3cm}
\textbf{Calculate 95\% CI:}

Standard error: $\text{SE} = \frac{8.4}{\sqrt{16}} = \frac{8.4}{4} = 2.10$

Critical value: $t_{0.975, 15} = 2.131$

Margin of error: $\text{ME} = 2.131 \times 2.10 = 4.48$

\textbf{95\% CI: $(45.2 - 4.48, 45.2 + 4.48) = (40.7, 49.7)$ points}

\vspace{0.3cm}
\textbf{Interpretation:} We are 95\% confident that the true mean pain reduction is between 40.7 and 49.7 points.
\end{frame}

\section{Factors Affecting CI Width}

\begin{frame}
\frametitle{What Makes a CI Wider or Narrower?}
\textbf{Recall:} $\text{CI} = \bar{x} \pm t_{\alpha/2, n-1} \times \frac{s}{\sqrt{n}}$

\vspace{0.5cm}
\textbf{Three main factors affect width:}

\begin{enumerate}
\item \textbf{Sample size ($n$)}
  \begin{itemize}
  \item Larger $n$ $\Rightarrow$ smaller SE $\Rightarrow$ narrower CI
  \item Width decreases by factor of $1/\sqrt{n}$
  \item To halve width, need $4\times$ the sample size
  \end{itemize}

\item \textbf{Sample variability ($s$)}
  \begin{itemize}
  \item Larger $s$ $\Rightarrow$ larger SE $\Rightarrow$ wider CI
  \item More variable data = more uncertainty
  \item Cannot control $s$ (it's a property of the population)
  \end{itemize}

\item \textbf{Confidence level ($1-\alpha$)}
  \begin{itemize}
  \item Higher confidence $\Rightarrow$ larger $t$ $\Rightarrow$ wider CI
  \item Trade-off: confidence vs. precision
  \end{itemize}
\end{enumerate}
\end{frame}

\begin{frame}
\frametitle{Effect of Sample Size}
\textbf{Same data, different sample sizes:}

\vspace{0.3cm}
Suppose $\bar{x} = 100$, $s = 20$ across all samples

\vspace{0.3cm}
\begin{center}
\begin{tabular}{cccccc}
\toprule
$n$ & SE & df & $t_{0.975}$ & ME & 95\% CI \\
\midrule
10 & 6.32 & 9 & 2.262 & 14.3 & (85.7, 114.3) \\
25 & 4.00 & 24 & 2.064 & 8.3 & (91.7, 108.3) \\
50 & 2.83 & 49 & 2.010 & 5.7 & (94.3, 105.7) \\
100 & 2.00 & 99 & 1.984 & 4.0 & (96.0, 104.0) \\
\bottomrule
\end{tabular}
\end{center}

\vspace{0.3cm}
\textbf{Observations:}
\begin{itemize}
\item As $n$ increases, SE decreases
\item As $n$ increases, $t$ approaches 1.96 (Z-value)
\item CI width decreases substantially with larger $n$
\item Diminishing returns: 10$\to$25 gives bigger gain than 50$\to$100
\end{itemize}
\end{frame}

\begin{frame}
\frametitle{Effect of Confidence Level}
\textbf{Same data, different confidence levels:}

\vspace{0.3cm}
Pain study: $n = 16$, $\bar{x} = 45.2$, $s = 8.4$, $\text{SE} = 2.10$

\vspace{0.3cm}
\begin{center}
\begin{tabular}{ccccc}
\toprule
Confidence & $\alpha$ & $t_{\alpha/2, 15}$ & ME & CI \\
\midrule
90\% & 0.10 & 1.753 & 3.68 & (41.5, 48.9) \\
95\% & 0.05 & 2.131 & 4.48 & (40.7, 49.7) \\
99\% & 0.01 & 2.947 & 6.19 & (39.0, 51.4) \\
\bottomrule
\end{tabular}
\end{center}

\vspace{0.5cm}
\textbf{The confidence-precision trade-off:}
\begin{itemize}
\item More confident $\Rightarrow$ wider interval (less precise)
\item Less confident $\Rightarrow$ narrower interval (more precise)
\item 95\% is conventional in most fields
\item 99\% for high-stakes decisions (e.g., safety-critical applications)
\item 90\% for preliminary or exploratory studies
\end{itemize}
\end{frame}

\begin{frame}
\frametitle{Effect of Variability}
\textbf{Same sample size and confidence, different variability:}

\vspace{0.3cm}
All samples: $n = 25$, $\bar{x} = 50$, 95\% confidence, $t_{0.975, 24} = 2.064$

\vspace{0.3cm}
\begin{center}
\begin{tabular}{ccccc}
\toprule
Scenario & $s$ & SE & ME & 95\% CI \\
\midrule
Low variability & 5 & 1.00 & 2.06 & (48.0, 52.0) \\
Moderate variability & 10 & 2.00 & 4.13 & (45.9, 54.1) \\
High variability & 20 & 4.00 & 8.26 & (41.7, 58.3) \\
\bottomrule
\end{tabular}
\end{center}

\vspace{0.5cm}
\textbf{Key insight:}
\begin{itemize}
\item Variability directly impacts precision
\item Cannot change $s$ through study design (it's a population property)
\item Can only counteract high $s$ by increasing $n$
\item More homogeneous populations give more precise estimates
\end{itemize}
\end{frame}

\section{Confidence Intervals for Proportions}

\begin{frame}
\frametitle{CI for a Population Proportion}
\textbf{When the outcome is binary:}
\begin{itemize}
\item Success/failure, yes/no, disease/no disease
\item Estimate population proportion $p$
\item Sample proportion: $\hat{p} = \frac{x}{n}$ (where $x$ = number of successes)
\end{itemize}

\vspace{0.5cm}
\textbf{Standard error for proportion:}
$$\text{SE} = \sqrt{\frac{\hat{p}(1-\hat{p})}{n}}$$

\vspace{0.5cm}
\textbf{95\% Confidence interval (Wald method):}
$$\hat{p} \pm z_{0.975} \times \sqrt{\frac{\hat{p}(1-\hat{p})}{n}}$$
$$\hat{p} \pm 1.96 \times \sqrt{\frac{\hat{p}(1-\hat{p})}{n}}$$

\vspace{0.3cm}
\textbf{Note:} Uses Z (not t) because with large $n$, proportion is approximately normal
\end{frame}

\begin{frame}
\frametitle{When to Use Normal Approximation for Proportions}
\textbf{The normal approximation works well when:}
$$n\hat{p} \geq 10 \quad \text{and} \quad n(1-\hat{p}) \geq 10$$

\vspace{0.3cm}
This ensures:
\begin{itemize}
\item At least 10 successes
\item At least 10 failures
\item Sampling distribution of $\hat{p}$ is approximately normal
\end{itemize}

\vspace{0.5cm}
\textbf{If conditions not met:}
\begin{itemize}
\item Use exact binomial methods (beyond this course)
\item Or use Wilson score interval (more advanced)
\item Or collect more data
\end{itemize}

\vspace{0.5cm}
\textbf{Common in biostatistics:}
\begin{itemize}
\item Disease prevalence
\item Treatment success rates
\item Adverse event rates
\item Survey response proportions
\end{itemize}
\end{frame}

\begin{frame}
\frametitle{Example: Vaccine Efficacy}
\textbf{Clinical trial:} New vaccine tested on 200 people
\begin{itemize}
\item 174 developed immunity (positive antibody response)
\item Sample proportion: $\hat{p} = \frac{174}{200} = 0.87$
\end{itemize}

\vspace{0.3cm}
\textbf{Check conditions:}
\begin{itemize}
\item $n\hat{p} = 200 \times 0.87 = 174 \geq 10$ \checkmark
\item $n(1-\hat{p}) = 200 \times 0.13 = 26 \geq 10$ \checkmark
\end{itemize}

\vspace{0.3cm}
\textbf{Calculate 95\% CI:}

Standard error: $\text{SE} = \sqrt{\frac{0.87 \times 0.13}{200}} = \sqrt{\frac{0.1131}{200}} = \sqrt{0.0005655} = 0.0238$

Margin of error: $\text{ME} = 1.96 \times 0.0238 = 0.047$

\textbf{95\% CI: $0.87 \pm 0.047 = (0.823, 0.917)$ or (82.3\%, 91.7\%)}

\vspace{0.3cm}
\textbf{Interpretation:} We are 95\% confident that the true vaccine efficacy is between 82.3\% and 91.7\%.
\end{frame}

\begin{frame}
\frametitle{Example: Disease Prevalence}
\textbf{Public health survey:} Screen 500 adults for diabetes
\begin{itemize}
\item 65 test positive
\item Sample proportion: $\hat{p} = \frac{65}{500} = 0.13$
\end{itemize}

\vspace{0.3cm}
\textbf{Check conditions:}
\begin{itemize}
\item $n\hat{p} = 500 \times 0.13 = 65 \geq 10$ \checkmark
\item $n(1-\hat{p}) = 500 \times 0.87 = 435 \geq 10$ \checkmark
\end{itemize}

\vspace{0.3cm}
\textbf{Calculate 95\% CI:}

SE: $\sqrt{\frac{0.13 \times 0.87}{500}} = \sqrt{0.0002262} = 0.0150$

ME: $1.96 \times 0.0150 = 0.0294$

\textbf{95\% CI: $0.13 \pm 0.029 = (0.101, 0.159)$ or (10.1\%, 15.9\%)}

\vspace{0.3cm}
\textbf{Interpretation:} We estimate that between 10.1\% and 15.9\% of adults in this population have diabetes.
\end{frame}

\section{Sample Size Determination}

\begin{frame}
\frametitle{Planning Studies: How Large a Sample?}
\textbf{Before collecting data, researchers ask:}
\begin{itemize}
\item How many subjects do we need?
\item What precision do we want?
\item How wide will our confidence interval be?
\end{itemize}

\vspace{0.5cm}
\textbf{Sample size determination:}
\begin{itemize}
\item Work backwards from desired margin of error
\item Requires estimate of population variability
\item Balances precision with cost/feasibility
\end{itemize}

\vspace{0.5cm}
\textbf{General approach:}
\begin{enumerate}
\item Specify desired margin of error (ME)
\item Estimate population SD ($\sigma$) from pilot data or literature
\item Choose confidence level (usually 95\%)
\item Solve for $n$ in the ME equation
\end{enumerate}
\end{frame}

\begin{frame}
\frametitle{Sample Size for Estimating a Mean}
\textbf{Starting with margin of error formula:}
$$\text{ME} = z_{\alpha/2} \times \frac{\sigma}{\sqrt{n}}$$

\vspace{0.3cm}
\textbf{Solve for $n$:}
$$\sqrt{n} = \frac{z_{\alpha/2} \times \sigma}{\text{ME}}$$
$$n = \left(\frac{z_{\alpha/2} \times \sigma}{\text{ME}}\right)^2$$

\vspace{0.5cm}
\textbf{For 95\% confidence:}
$$n = \left(\frac{1.96 \times \sigma}{\text{ME}}\right)^2$$

\vspace{0.3cm}
\textbf{Note:}
\begin{itemize}
\item Uses Z (not t) for planning
\item Round up to next integer
\item May need to iterate if using t-distribution
\end{itemize}
\end{frame}

\begin{frame}
\frametitle{Example: Cholesterol Study Planning}
\textbf{Goal:} Estimate mean cholesterol with margin of error = 5 mg/dL

\vspace{0.3cm}
\textbf{Given information:}
\begin{itemize}
\item Desired ME = 5 mg/dL
\item From literature: $\sigma \approx 40$ mg/dL
\item Confidence level: 95\% (so $z_{0.975} = 1.96$)
\end{itemize}

\vspace{0.3cm}
\textbf{Calculate required sample size:}
$$n = \left(\frac{1.96 \times 40}{5}\right)^2 = \left(\frac{78.4}{5}\right)^2 = (15.68)^2 = 245.86$$

\vspace{0.3cm}
\textbf{Round up:} $n = 246$ subjects

\vspace{0.5cm}
\textbf{Interpretation:}
\begin{itemize}
\item Need to sample 246 people
\item Will give 95\% CI with margin of error $\pm 5$ mg/dL
\item Actual CI might be slightly different (using $s$ instead of $\sigma$)
\end{itemize}
\end{frame}

\begin{frame}
\frametitle{Example: Blood Pressure Study Planning}
\textbf{Different margins of error:}

\vspace{0.3cm}
Suppose $\sigma = 15$ mmHg for blood pressure, 95\% confidence

\vspace{0.3cm}
\begin{center}
\begin{tabular}{ccc}
\toprule
Desired ME & Calculation & Required $n$ \\
\midrule
2 mmHg & $(1.96 \times 15 / 2)^2$ & 217 \\
3 mmHg & $(1.96 \times 15 / 3)^2$ & 97 \\
5 mmHg & $(1.96 \times 15 / 5)^2$ & 35 \\
10 mmHg & $(1.96 \times 15 / 10)^2$ & 9 \\
\bottomrule
\end{tabular}
\end{center}

\vspace{0.5cm}
\textbf{Key insights:}
\begin{itemize}
\item Halving ME requires $4\times$ the sample size
\item More precision is expensive (in terms of sample size)
\item Diminishing returns: going from ME=5 to ME=2 adds 182 subjects
\item Practical considerations: balance precision with feasibility
\end{itemize}
\end{frame}

\begin{frame}
\frametitle{Sample Size for Estimating a Proportion}
\textbf{For proportions:}
$$n = \left(\frac{z_{\alpha/2}}{\text{ME}}\right)^2 p(1-p)$$

\vspace{0.3cm}
\textbf{Problem:} We need $p$ to calculate $n$, but we don't know $p$ yet!

\vspace{0.5cm}
\textbf{Three approaches:}

\begin{enumerate}
\item \textbf{Use pilot data or prior estimate:}
  \begin{itemize}
  \item Previous studies suggest $p \approx 0.3$
  \item Use this value in the formula
  \end{itemize}

\item \textbf{Conservative approach:}
  \begin{itemize}
  \item Use $p = 0.5$ (maximizes $p(1-p)$)
  \item Gives largest possible $n$ (safe but possibly wasteful)
  \end{itemize}

\item \textbf{Range of plausible values:}
  \begin{itemize}
  \item Calculate $n$ for several plausible $p$ values
  \item Choose based on context
  \end{itemize}
\end{enumerate}
\end{frame}

\begin{frame}
\frametitle{Example: Prevalence Study Planning}
\textbf{Goal:} Estimate disease prevalence with ME = 0.03 (3\%), 95\% confidence

\vspace{0.3cm}
\textbf{Scenario 1:} No prior information (use $p = 0.5$)
$$n = \left(\frac{1.96}{0.03}\right)^2 \times 0.5 \times 0.5 = (65.33)^2 \times 0.25 = 1067.1$$
Need $n = 1068$ subjects

\vspace{0.3cm}
\textbf{Scenario 2:} Prior studies suggest $p \approx 0.15$
$$n = \left(\frac{1.96}{0.03}\right)^2 \times 0.15 \times 0.85 = (65.33)^2 \times 0.1275 = 544.5$$
Need $n = 545$ subjects

\vspace{0.5cm}
\textbf{Comparison:}
\begin{itemize}
\item Conservative approach requires nearly 2$\times$ more subjects
\item If we have good prior information, can save resources
\item Trade-off: certainty in planning vs. sample size efficiency
\end{itemize}
\end{frame}

\section{Advanced Topics and Considerations}

\begin{frame}
\frametitle{Large Sample vs Small Sample CIs}
\textbf{When $n$ is large ($n \geq 30$ as rule of thumb):}
\begin{itemize}
\item t-distribution $\approx$ normal distribution
\item Can use $z \approx 1.96$ instead of $t$ for 95\% CI
\item CLT ensures $\bar{X}$ is approximately normal
\item Less sensitive to departures from normality
\end{itemize}

\vspace{0.5cm}
\textbf{When $n$ is small ($n < 30$):}
\begin{itemize}
\item Must use t-distribution (heavier tails)
\item Need stronger normality assumption
\item Check data for severe skewness or outliers
\item Consider transformations if badly non-normal
\end{itemize}

\vspace{0.5cm}
\textbf{Practical guideline:}
\begin{itemize}
\item Always use t-distribution when $\sigma$ unknown
\item Check normality assumption with histogram or Q-Q plot
\item For very small $n$ ($< 15$), be especially careful
\item Robust methods exist for non-normal data (beyond this course)
\end{itemize}
\end{frame}

\begin{frame}
\frametitle{Using CIs for Hypothesis Testing}
\textbf{CIs can answer yes/no questions:}

\vspace{0.3cm}
\textbf{Example:} Is the population mean different from a specific value?

\vspace{0.3cm}
\textbf{Drug efficacy study:}
\begin{itemize}
\item Research question: Does the drug reduce pain by more than 30 points on average?
\item 95\% CI for mean reduction: (35.2, 48.6)
\item Since entire CI is above 30, YES, strong evidence
\end{itemize}

\vspace{0.5cm}
\textbf{General rule for 95\% CI:}
\begin{itemize}
\item If hypothesized value is \textit{outside} the CI: reject at $\alpha = 0.05$
\item If hypothesized value is \textit{inside} the CI: fail to reject at $\alpha = 0.05$
\item CI gives more information than p-value (shows magnitude and precision)
\end{itemize}

\vspace{0.3cm}
\textbf{Next week:} Formal hypothesis testing procedures
\end{frame}

\begin{frame}
\frametitle{Common Mistakes with Confidence Intervals}
\textbf{1. Misinterpreting the confidence level}
\begin{itemize}
\item $\times$ ``95\% probability that $\mu$ is in this interval''
\item \checkmark ``95\% of such intervals contain $\mu$''
\end{itemize}

\vspace{0.3cm}
\textbf{2. Confusing CI for mean with range of data}
\begin{itemize}
\item CI estimates where population mean is
\item Not where individual observations fall
\end{itemize}

\vspace{0.3cm}
\textbf{3. Ignoring assumptions}
\begin{itemize}
\item Random sampling required
\item Independence required
\item Approximate normality (or large $n$)
\end{itemize}

\vspace{0.3cm}
\textbf{4. Using Z when should use t}
\begin{itemize}
\item Always use t when $\sigma$ is unknown (almost always!)
\end{itemize}

\vspace{0.3cm}
\textbf{5. Forgetting to round $n$ up in planning}
\begin{itemize}
\item $n = 245.86$ means need 246 subjects, not 245
\end{itemize}
\end{frame}

\begin{frame}
\frametitle{Checking Assumptions: Normality}
\textbf{How to assess if data are approximately normal:}

\vspace{0.3cm}
\textbf{1. Histogram}
\begin{itemize}
\item Should be roughly bell-shaped
\item Check for severe skewness or multiple modes
\end{itemize}

\vspace{0.3cm}
\textbf{2. Q-Q plot (quantile-quantile plot)}
\begin{itemize}
\item Points should fall roughly on a straight line
\item Deviations at ends less concerning than in middle
\end{itemize}

\vspace{0.3cm}
\textbf{3. Formal tests (Shapiro-Wilk)}
\begin{itemize}
\item Tests null hypothesis that data are normal
\item Very sensitive with large $n$ (may reject for trivial departures)
\end{itemize}

\vspace{0.3cm}
\textbf{Remember:}
\begin{itemize}
\item t-procedures are robust to mild departures
\item Larger $n$ means less sensitive to non-normality (CLT)
\item Outliers more problematic than skewness
\end{itemize}
\end{frame}

\section{AI as a Learning Tool for Confidence Intervals}

\AIcolors
\begin{frame}
\frametitle{AI for Confidence Interval Analysis}
\textbf{How AI can help this week:}
\begin{itemize}
\item \textbf{CI calculations:} Computing intervals step-by-step
\item \textbf{Sample size planning:} Determining required $n$
\item \textbf{Interpretation:} Writing correct interpretations
\item \textbf{Assumption checking:} Creating diagnostic plots
\item \textbf{Comparison:} Different confidence levels or methods
\item \textbf{R code:} Using \texttt{t.test()} and related functions
\end{itemize}

\vspace{0.3cm}
\textbf{Example prompts:}
\begin{itemize}
\item ``Calculate 95\% CI for mean given: n=20, mean=105, SD=15''
\item ``How many subjects needed for ME=3 if SD=12 at 95\% confidence?''
\item ``Create R code to check normality assumption with histogram and Q-Q plot''
\item ``Explain the difference between 90\%, 95\%, and 99\% confidence intervals''
\end{itemize}
\end{frame}

\begin{frame}[fragile]
\frametitle{AI Demo: CI Calculation in R}
\textbf{Prompt:} ``Calculate 95\% CI for cholesterol data in R''

\vspace{0.2cm}
\textbf{Claude AI Response:}
\tiny
\begin{verbatim}
# Cholesterol data (mg/dL)
cholesterol <- c(205, 198, 220, 185, 210, 195, 225, 200, 215, 192,
                 208, 196, 218, 203, 212, 189, 207, 201, 213, 199)

# Calculate sample statistics
n <- length(cholesterol)
xbar <- mean(cholesterol)
s <- sd(cholesterol)
se <- s / sqrt(n)

# Calculate 95% CI using t-distribution
df <- n - 1
t_crit <- qt(0.975, df)
me <- t_crit * se
ci_lower <- xbar - me
ci_upper <- xbar + me

# Display results
cat("Sample size:", n, "\n")
cat("Sample mean:", round(xbar, 2), "mg/dL\n")
cat("Sample SD:", round(s, 2), "mg/dL\n")
cat("Standard error:", round(se, 3), "mg/dL\n")
cat("t critical value:", round(t_crit, 3), "\n")
cat("Margin of error:", round(me, 2), "mg/dL\n")
cat("95% CI: (", round(ci_lower, 2), ", ", round(ci_upper, 2), ")\n", sep="")

# Alternative: use t.test() function
t.test(cholesterol, conf.level=0.95)
\end{verbatim}
\end{frame}

\begin{frame}[fragile]
\frametitle{AI Demo: Sample Size Calculation}
\textbf{Prompt:} ``Calculate sample size needed for blood pressure study in R''

\vspace{0.2cm}
\textbf{Claude AI Response:}
\tiny
\begin{verbatim}
# Sample size calculation for estimating mean BP
# Goal: 95% CI with margin of error = 3 mmHg
# Estimated population SD = 15 mmHg

sigma_est <- 15     # Estimated population SD
me_desired <- 3     # Desired margin of error
conf_level <- 0.95  # Confidence level
alpha <- 1 - conf_level
z_crit <- qnorm(1 - alpha/2)  # 1.96 for 95%

# Calculate required sample size
n_required <- (z_crit * sigma_est / me_desired)^2

cat("Sample size calculation:\n")
cat("Estimated SD:", sigma_est, "mmHg\n")
cat("Desired margin of error:", me_desired, "mmHg\n")
cat("Confidence level:", conf_level*100, "%\n")
cat("Z critical value:", round(z_crit, 3), "\n\n")
cat("Required sample size (exact):", round(n_required, 2), "\n")
cat("Required sample size (rounded up):", ceiling(n_required), "\n")

# Verify the calculation
actual_me <- z_crit * sigma_est / sqrt(ceiling(n_required))
cat("\nVerification - actual ME with n =", ceiling(n_required), ":",
    round(actual_me, 3), "mmHg\n")
\end{verbatim}
\end{frame}

\begin{frame}[fragile]
\frametitle{AI Demo: Comparing Confidence Levels}
\textbf{Prompt:} ``Compare 90\%, 95\%, 99\% CIs for same data in R''

\vspace{0.2cm}
\textbf{Claude AI Response:}
\tiny
\begin{verbatim}
# Sample data: pain reduction scores
pain_reduction <- c(42, 38, 51, 45, 39, 48, 44, 50, 43, 47,
                    41, 46, 49, 40, 52)

# Calculate sample statistics
n <- length(pain_reduction)
xbar <- mean(pain_reduction)
s <- sd(pain_reduction)
se <- s / sqrt(n)
df <- n - 1

# Calculate CIs for different confidence levels
conf_levels <- c(0.90, 0.95, 0.99)

cat("Sample mean:", round(xbar, 2), "\n")
cat("Sample SD:", round(s, 3), "\n")
cat("Sample size:", n, "\n")
cat("Standard error:", round(se, 3), "\n\n")

for (conf in conf_levels) {
  alpha <- 1 - conf
  t_crit <- qt(1 - alpha/2, df)
  me <- t_crit * se
  ci_lower <- xbar - me
  ci_upper <- xbar + me
  width <- ci_upper - ci_lower

  cat(conf*100, "% CI: (", round(ci_lower, 2), ", ",
      round(ci_upper, 2), ")  Width = ", round(width, 2), "\n", sep="")
}

cat("\nNotice: Higher confidence = wider interval\n")
\end{verbatim}
\end{frame}

\begin{frame}[fragile]
\frametitle{AI Demo: Checking Normality Assumption}
\textbf{Prompt:} ``Create diagnostic plots to check normality assumption''

\vspace{0.2cm}
\textbf{Claude AI Response:}
\tiny
\begin{verbatim}
# Sample data
data <- c(23, 25, 28, 22, 27, 24, 26, 29, 25, 23, 28, 24, 27, 25, 26)

# Create diagnostic plots
par(mfrow=c(1,2))  # 1 row, 2 columns

# 1. Histogram with normal curve overlay
hist(data, breaks=8, prob=TRUE, col="lightblue",
     main="Histogram with Normal Curve",
     xlab="Value", border="white")
# Overlay theoretical normal curve
xfit <- seq(min(data), max(data), length=100)
yfit <- dnorm(xfit, mean=mean(data), sd=sd(data))
lines(xfit, yfit, col="red", lwd=2)

# 2. Q-Q plot
qqnorm(data, main="Q-Q Plot", pch=19, col="blue")
qqline(data, col="red", lwd=2)

par(mfrow=c(1,1))  # Reset layout

# Formal normality test
shapiro_test <- shapiro.test(data)
cat("\nShapiro-Wilk test for normality:\n")
cat("W statistic:", round(shapiro_test$statistic, 4), "\n")
cat("p-value:", round(shapiro_test$p.value, 4), "\n")
if (shapiro_test$p.value > 0.05) {
  cat("Conclusion: No evidence against normality (p > 0.05)\n")
} else {
  cat("Conclusion: Evidence of non-normality (p < 0.05)\n")
}
\end{verbatim}
\end{frame}

\begin{frame}
\frametitle{AI Best Practices for CIs}
\begin{columns}
\begin{column}{0.5\textwidth}
\textbf{DO:}
\begin{itemize}
\item Verify calculations manually first
\item Check assumptions before computing CI
\item Ask for interpretations
\item Request visualizations
\item Compare different methods
\item Use AI to learn concepts
\end{itemize}
\end{column}

\begin{column}{0.5\textwidth}
\textbf{DON'T:}
\begin{itemize}
\item Trust results blindly
\item Skip assumption checking
\item Confuse Z and t methods
\item Ignore sample size requirements
\item Forget to interpret in context
\item Use without understanding
\end{itemize}
\end{column}
\end{columns}

\vspace{0.3cm}
\textbf{Key Verification Steps:}
\begin{enumerate}
\item Check that sample statistics are correct
\item Verify correct critical value (t not Z, correct df)
\item Confirm SE calculation
\item Ensure interpretation makes sense in context
\item Cross-check with \texttt{t.test()} output
\end{enumerate}
\end{frame}

\defaultcolors

\section{R Examples and Practice}

\begin{frame}[fragile]
\frametitle{R Example: Basic t-Interval}
\small
\begin{verbatim}
# Blood pressure data (mmHg)
bp <- c(128, 132, 125, 138, 130, 135, 127, 140,
        129, 133, 136, 126, 131, 134, 137)

# Method 1: Manual calculation
xbar <- mean(bp)
s <- sd(bp)
n <- length(bp)
se <- s / sqrt(n)
df <- n - 1
t_crit <- qt(0.975, df)  # 95% CI
me <- t_crit * se
ci <- c(xbar - me, xbar + me)

cat("Manual: 95% CI =", round(ci, 2), "\n")

# Method 2: Using t.test()
result <- t.test(bp, conf.level=0.95)
cat("t.test(): 95% CI =", round(result$conf.int, 2), "\n")

# Both methods give same answer
print(result)
\end{verbatim}
\end{frame}

\begin{frame}[fragile]
\frametitle{R Example: CI for Proportion}
\small
\begin{verbatim}
# Vaccine study: 174 out of 200 developed immunity
x <- 174        # Number of successes
n <- 200        # Sample size
p_hat <- x / n  # Sample proportion

# Check conditions
cat("np =", n*p_hat, ", n(1-p) =", n*(1-p_hat), "\n")

# Calculate 95% CI
se <- sqrt(p_hat * (1 - p_hat) / n)
z_crit <- qnorm(0.975)  # 1.96
me <- z_crit * se
ci_lower <- p_hat - me
ci_upper <- p_hat + me

cat("Sample proportion:", p_hat, "\n")
cat("Standard error:", round(se, 4), "\n")
cat("95% CI: (", round(ci_lower, 4), ", ",
    round(ci_upper, 4), ")\n", sep="")
cat("95% CI: (", round(ci_lower*100, 2), "%, ",
    round(ci_upper*100, 2), "%)\n", sep="")

# Alternative: use prop.test()
prop.test(x, n, conf.level=0.95)
\end{verbatim}
\end{frame}

\begin{frame}[fragile]
\frametitle{R Example: Sample Size Planning}
\footnotesize
\begin{verbatim}
# Function to calculate required sample size for mean
sample_size_mean <- function(sigma, me, conf.level=0.95) {
  alpha <- 1 - conf.level
  z <- qnorm(1 - alpha/2)
  n <- (z * sigma / me)^2
  return(ceiling(n))  # Round up
}

# Example: Cholesterol study
sigma_est <- 40   # mg/dL
me_desired <- 5   # mg/dL

n_needed <- sample_size_mean(sigma_est, me_desired, 0.95)
cat("Required sample size:", n_needed, "\n")

# Try different margins of error
margins <- c(2, 3, 5, 10)
for (me in margins) {
  n <- sample_size_mean(sigma_est, me, 0.95)
  cat("ME =", me, "-> n =", n, "\n")
}

# For proportions
sample_size_prop <- function(p, me, conf.level=0.95) {
  alpha <- 1 - conf.level
  z <- qnorm(1 - alpha/2)
  n <- (z / me)^2 * p * (1-p)
  return(ceiling(n))
}

# Conservative (p=0.5)
cat("\nProportion (p=0.5, ME=0.03):",
    sample_size_prop(0.5, 0.03, 0.95), "\n")
\end{verbatim}
\end{frame}

\begin{frame}[fragile]
\frametitle{R Example: Visualizing Multiple CIs}
\footnotesize
\begin{verbatim}
# Simulate multiple samples and CIs
set.seed(524)
true_mean <- 100   # True population mean
true_sd <- 15      # True population SD
n <- 25            # Sample size
num_samples <- 20  # Number of samples to draw

# Storage for CIs
ci_data <- data.frame(sample=1:num_samples,
                      lower=NA, upper=NA, contains_mu=NA)

for (i in 1:num_samples) {
  sample_data <- rnorm(n, true_mean, true_sd)
  result <- t.test(sample_data, conf.level=0.95)
  ci_data$lower[i] <- result$conf.int[1]
  ci_data$upper[i] <- result$conf.int[2]
  ci_data$contains_mu[i] <- (ci_data$lower[i] <= true_mean) &
                             (true_mean <= ci_data$upper[i])
}

# Plot the CIs
plot(NULL, xlim=c(85, 115), ylim=c(0, num_samples+1),
     xlab="Value", ylab="Sample Number",
     main="20 Confidence Intervals (95% level)")
abline(v=true_mean, col="blue", lwd=2)
for (i in 1:num_samples) {
  color <- ifelse(ci_data$contains_mu[i], "darkgreen", "red")
  segments(ci_data$lower[i], i, ci_data$upper[i], i, col=color, lwd=2)
  points(mean(c(ci_data$lower[i], ci_data$upper[i])), i, pch=19, col=color)
}
legend("topleft", c("True mean", "Contains mu", "Misses mu"),
       col=c("blue", "darkgreen", "red"), lty=1, lwd=2)
cat(sum(ci_data$contains_mu), "out of", num_samples, "CIs contain mu\n")
\end{verbatim}
\end{frame}

\begin{frame}
\frametitle{Summary}
\textbf{Chapter 9 Key Concepts:}
\begin{itemize}
\item \textbf{Confidence intervals:} Range of plausible values for parameter
  \begin{itemize}
  \item General form: Estimate $\pm$ (Critical value) $\times$ SE
  \end{itemize}
\item \textbf{t-interval for mean:} $\bar{x} \pm t_{\alpha/2, n-1} \times \frac{s}{\sqrt{n}}$
  \begin{itemize}
  \item Use when $\sigma$ unknown (almost always)
  \item Requires approximate normality or large $n$
  \end{itemize}
\item \textbf{CI for proportion:} $\hat{p} \pm z_{\alpha/2} \times \sqrt{\frac{\hat{p}(1-\hat{p})}{n}}$
  \begin{itemize}
  \item Requires $n\hat{p} \geq 10$ and $n(1-\hat{p}) \geq 10$
  \end{itemize}
\item \textbf{Factors affecting width:}
  \begin{itemize}
  \item Sample size $n$ (larger $\to$ narrower)
  \item Variability $s$ (larger $\to$ wider)
  \item Confidence level (higher $\to$ wider)
  \end{itemize}
\item \textbf{Sample size planning:} Work backwards from desired ME
\end{itemize}

\vspace{0.3cm}
\textbf{Next Week:} Hypothesis testing (one-sample and two-sample tests)
\end{frame}

\begin{frame}
\frametitle{Questions?}
\begin{center}
\Large{Thank you!}

\vspace{1cm}

\textbf{Email:} john.molitor@oregonstate.edu\\
\textbf{Office:} 153 Milam
\end{center}
\end{frame}

\end{document}
